Table~\ref{tab:main} gives the main results (mean $\pm$ std over three seeds) and Fig.~\ref{fig:main} the held-out accuracy. Significance tests below are Welch tests from \texttt{rh compare} with $n=3$ per cell; with so few seeds they are weak evidence.

\begin{table*}[t]\centering\caption{Main results, mean $\pm$ std over 3 seeds. Oracle and random code are fixed senders (references), not competitors.}\label{tab:main}
\resizebox{0.85\textwidth}{!}{\begin{tabular}{lcccc}
\toprule
Task / system & held-out & train & topsim & gap \\
\midrule
V4\_L4 Gumbel & 0.013 $\pm$ 0.022 & 0.052 $\pm$ 0.033 & 0.101 $\pm$ 0.090 & 0.039 $\pm$ 0.042 \\
V4\_L4 REINFORCE & 0.013 $\pm$ 0.022 & 0.433 $\pm$ 0.108 & 0.300 $\pm$ 0.018 & 0.421 $\pm$ 0.104 \\
V4\_L4 Random code & 0.000 $\pm$ 0.000 & 0.530 $\pm$ 0.012 & 0.004 $\pm$ 0.002 & 0.530 $\pm$ 0.012 \\
V4\_L4 Oracle code & 1.000 $\pm$ 0.000 & 1.000 $\pm$ 0.000 & 1.000 $\pm$ 0.000 & 0.000 $\pm$ 0.000 \\
V8\_L4 Gumbel & 0.026 $\pm$ 0.022 & 0.381 $\pm$ 0.147 & 0.236 $\pm$ 0.060 & 0.356 $\pm$ 0.169 \\
V8\_L4 REINFORCE & 0.000 $\pm$ 0.000 & 0.449 $\pm$ 0.124 & 0.296 $\pm$ 0.045 & 0.449 $\pm$ 0.124 \\
V8\_L4 Random code & 0.000 $\pm$ 0.000 & 0.968 $\pm$ 0.015 & -0.000 $\pm$ 0.004 & 0.968 $\pm$ 0.015 \\
V8\_L4 Oracle code & 1.000 $\pm$ 0.000 & 1.000 $\pm$ 0.000 & 1.000 $\pm$ 0.000 & 0.000 $\pm$ 0.000 \\
V16\_L8 Gumbel & 0.436 $\pm$ 0.232 & 0.775 $\pm$ 0.135 & 0.264 $\pm$ 0.020 & 0.339 $\pm$ 0.179 \\
V16\_L8 REINFORCE & 0.500 $\pm$ 0.176 & 0.991 $\pm$ 0.009 & 0.364 $\pm$ 0.012 & 0.491 $\pm$ 0.171 \\
V16\_L8 Random code & 0.000 $\pm$ 0.000 & 1.000 $\pm$ 0.000 & 0.004 $\pm$ 0.005 & 1.000 $\pm$ 0.000 \\
V16\_L8 Oracle code & 0.846 $\pm$ 0.038 & 1.000 $\pm$ 0.000 & 1.000 $\pm$ 0.000 & 0.154 $\pm$ 0.038 \\
\bottomrule
\end{tabular}
}\end{table*}

\begin{figure}[t]\centering\includegraphics[width=\columnwidth]{main_bars.pdf}
\caption{Held-out accuracy (left) and topsim (right) by system and task, mean and std over 3 seeds.}\label{fig:main}\end{figure}

\textbf{H1 (capacity needed): partly supported, confounded by optimisation.} In the sweeps (Tables~\ref{tab:sv},\ref{tab:sl}), channels below 8 bits (V2\_L5, V3\_L5, V6\_L2, V6\_L3) have train accuracy at most 0.094 and held-out accuracy at most 0.013. But accuracy does not switch on at 8 bits: V4\_L5 (10 bits) still has train accuracy 0.126, and V6\_L4 (10.3 bits) 0.172. Accuracy rises only for much larger channels (V16\_L5: 0.949 train). So capacity is necessary, but at this budget it is far from sufficient; ``capacity below 8 bits caps accuracy'' is true but uninformative.

\textbf{H2 (larger channel hurts generalisation): refuted in this study, with a caveat.} V16\_L8 reaches held-out accuracy 0.436 (Gumbel-softmax) and 0.500 (REINFORCE) while the tight V4\_L4 reaches 0.013 for both. The train--held-out gap is smaller at V4\_L4 (0.039 for Gumbel-softmax) only because train accuracy is itself 0.052: the agents have not learned the task there. Topsim is also higher, not lower, at V16\_L8 (0.264) than at V4\_L4 (0.101) for Gumbel-softmax. The over-capacity channel does carry a large gap (0.339 and 0.491), and the random code shows the extreme: train accuracy 1.000 at V16\_L8 but held-out 0.000. The data do not tell us whether a well-trained tight channel would generalise better. The curves (Fig.~\ref{fig:curves}, Table~\ref{tab:curves}) show why: Gumbel-softmax at V4\_L4 is unstable rather than slow. Train accuracy peaks at 0.239--0.291 between steps 600 and 1400 and then collapses (0.017--0.083 at step 6000, and at most 0.009 at 20000, with one or two distinct messages in the control runs of Table~\ref{tab:long}); the 300-step sanity run (train 0.174) beat the 6000-step run of the same seed (0.087). Longer training therefore does not help this system under $\tau=1$, lr $2\times10^{-3}$; we did not test other temperatures or learning rates, so the cause (e.g.\ a collapsing sender or temperature) is not identified.

\textbf{H3 (learned protocols are not compositional): supported.} Topsim of learned protocols lies between 0.101 and 0.364, versus 1.000 for the oracle code (Welch $p<0.01$ for every task, Gumbel-softmax vs oracle). On V4\_L4 and V8\_L4 the oracle reaches 1.000 held-out accuracy against at most 0.026 for learned senders. However the oracle is \emph{not} perfect at V16\_L8 (held-out 0.846): a listener GRU reading 8 positions of which 4 are constant padding generalises imperfectly, a failure of the reference listener rather than of the code that we did not investigate further. The random code has topsim near zero (0.004) and held-out accuracy 0.000 everywhere, as expected for an unstructured code.

\textbf{H4 (Gumbel-softmax beats REINFORCE): not supported.} On held-out accuracy the difference is within noise on every task (V16\_L8: 0.436 vs 0.500, Welch $p=0.72$). REINFORCE has clearly higher train accuracy on V4\_L4 (0.433 vs 0.052), where Gumbel-softmax collapses; at V8\_L4 the train-accuracy difference (0.449 vs 0.381) is within one standard deviation, and the topsim differences are within noise at V4\_L4 and V8\_L4 (Welch $p=0.057$ and $0.24$); only the V16\_L8 topsim difference (REINFORCE higher, $p=0.004$) is clear. Unlike Havrylov and Titov, who found Gumbel-softmax converging much faster, we did not observe an advantage for it. REINFORCE is still improving at 6000 steps at V4\_L4 (Table~\ref{tab:curves}), but its held-out accuracy stays at 0.000 after 20000 steps (Table~\ref{tab:long}) because it memorises (train 0.50) rather than generalises. With three seeds and learning rate, temperature and entropy weight untuned for either, we make no claim about the estimators in general.

\textbf{Failure cases.} Per-seed variability is large: the std of Gumbel-softmax held-out accuracy at V16\_L8 is 0.232 around a mean of 0.436, so a different seed set could change the ranking with REINFORCE. All learned systems at V4\_L4 and V8\_L4 fail (held-out at most 0.026), so no statement about compositionality at tight capacity can be drawn from well-trained protocols.

\begin{figure}[t]\centering\includegraphics[width=\columnwidth]{curves_v4l4.pdf}
\caption{Train accuracy (greedy messages) at V4\_L4 over 20000 steps, per seed, for Gumbel-softmax and REINFORCE.}\label{fig:curves}\end{figure}
\begin{table}[t]\centering\caption{Train-accuracy curves at V4\_L4: peak value and step, and value at steps 6000 and 20000 (from \texttt{results/raw/curve\_*.csv}).}\label{tab:curves}
\resizebox{\columnwidth}{!}{\begin{tabular}{llrrrr}
\hline
System & Seed & Peak & Peak step & @6000 & @20000\\
\hline
Gumbel & 0 & 0.265 & 1400 & 0.083 & 0.009\\
Gumbel & 1 & 0.291 & 1200 & 0.017 & 0.000\\
Gumbel & 2 & 0.239 & 600 & 0.043 & 0.009\\
REINFORCE & 0 & 0.357 & 16500 & 0.313 & 0.357\\
REINFORCE & 1 & 0.574 & 16600 & 0.526 & 0.570\\
REINFORCE & 2 & 0.570 & 19500 & 0.465 & 0.570\\
\hline
\end{tabular}
}\end{table}
\begin{table}[t]\centering\caption{Longer-training control at V4\_L4 (20000 steps, mean $\pm$ std over 3 seeds; msgs = distinct messages over the 256 objects).}\label{tab:long}
\resizebox{\columnwidth}{!}{\begin{tabular}{lrrrr}
\hline
System & train & held-out & topsim & msgs\\
\hline
Gumbel & 0.006 $\pm$ 0.005 & 0.013 $\pm$ 0.022 & 0.000 $\pm$ 0.000 & 1.7 $\pm$ 0.6\\
REINFORCE & 0.499 $\pm$ 0.123 & 0.000 $\pm$ 0.000 & 0.308 $\pm$ 0.024 & 117.0 $\pm$ 27.7\\
\hline
\end{tabular}
}\end{table}
